% Pursuing Stacks (cotes 134-2 to 134-8): archive cote 134-5, site page 46 (PDF page 47), typescript page 233.
% typescript. Draft compared: ps4.tex:2262-2327.
% AI-transcribed by Claude Opus 5.5 (claude-opus-5-5, adaptive thinking, effort high), Message Batches API, batch ps2 (msgbatch_01NVwr1Yuo8qJf5rykeGBXPK), 2026-11-12; not yet checked by eye.
\scanpage{134-5/46}
\mspage{233}
it's small subcategories should be in (Cat). This presumably will be sa-%
tisfied by most large categories we are going to consider, and it should\hand{\unsure{'|)}}
nt be hard moreoever to show that any $\underline{U}$-category is isomorphic to a cate-%
gory $M'$ satisfying the above extra assumption.

A little more serious maybe is the use we are making here of the
name of a ``testing functor'', which seems to be conflicting with an ear-%
lier use (def.5, p. 175\xout{)xxx} and def.6, p.176), where we insisted \inserted{(for instance)} that
the source $A$ should be a test category. That's why I have been using
here the name ``testing functor'' rather than ``test functor'', to be on
the safe side formally speaking - but this is still playing on words,
namely cheating a little. Maybe the name ``\uline{aspherical functor}'' instead
of``testing functor'' would be less misleading, thinking of the case when
$M$ is small itself, and endowing $M$ with the canonical asphericity struc-%
ture, for which
\[ M_{as} = M \quad , \]
(admitting the identity functor as a testing functor) - in which case
the``testing functors'' $A \to M$ are indeed just the aspherical functors.
The drawback is that when $M$ is small, \xout{there} it may well be endowed with
an asphericity structure $M_{as}$ different from the previous one, in which
case the proposed extension of the name ``aspheric functor'' again leads to
an ambiguity, unless specified by ``$M_{as}$-\uline{aspheric}''(where, after all,
$M_{as}$ could be \uline{any} full subcategory of $M$) - but then the notion reduces
to the one of an aspheric functor $A \to M_{as}$ . But the same after all
holds even for large $M$ - the notion of a ``testing functor'' $A \to M$
(with respect to an asphericity structure $M_{as}$ on $M$) does not really de-%
pend on the \uline{pair} $(M,M_{as})$, but rather on the (\xout{generally}\inserted[typed]{possibly} large) category
$M_{as}$ itself - namely it is no more no less than a functor
\[ A \to M_{as} \]
which satisfies the usual asphericity condition (iii) (of prop.1 p.214)\hand{,}
with the only difference that $M_{as}$ may not be small, and therefore \xout{we}
$M_{as}\hat{}$ may not be a $\underline{U}$-category (and we are \inserted{therefore} reluctant to work with this
latter category at all, unless we first pass to the next larger universe
$\underline{U}'$ ...).

This short reflection rather convinces me that the designation of
``testing functors'' as introduced on the page before, by the alternative
name of $M_{as}$-\uline{aspheric functors}, or just \uline{aspheric functors} when no confu-%
sion is likely to arise as to existence and choice of $M_{as}$, is satisfac-%
tory indeed. I'll use it tentatively, as a synonyme to ``testing func-%
tor'', and it will appear soon enough if this terminology is a good one
\margin{46}
